% =====================================================================
% research_summary.tex — Electoral Equilibrium
% ~2-page standalone project summary (prose). No fabricated numbers:
% real values or \todo{}. Compile: pdflatex research_summary.tex
% =====================================================================
\documentclass[11pt]{article}

\usepackage[margin=1in]{geometry}
\usepackage[dvipsnames]{xcolor}
\usepackage{amsmath,amssymb}
\usepackage{parskip}

\newcommand{\todo}[1]{{\color{red}\textbf{[TODO:} #1\textbf{]}}}

\title{Electoral Equilibrium: A Two-Page Summary}
\author{Juhi Damley, Claremont McKenna College; supervised by
Prof.\ Gaston Espinosa}
\date{\today}

\begin{document}
\maketitle

\section*{The problem}
Political shocks --- scandals, endorsements, economic surprises, institutional
rulings --- do not move the electorate uniformly; different demographic blocs
react differently. Electoral Equilibrium formalizes a party's most defensible
strategic coalition as an \emph{equilibrium}: the race-bloc weighting that
maximizes win probability given the party's baseline bloc loyalties, computed
once per party rather than once per shock. It then asks a single question:
given a hypothetical shock described in plain text, how does the shock's
predicted per-bloc loyalty shift move that fixed equilibrium coalition's
effective loyalty and win probability relative to its win threshold --- not
whether the coalition weighting itself should change in response, which, as
reported below, it empirically does not. Is the estimate partisan-neutral, and
how well does it match the magnitude and direction of real historical shifts?
The hard constraint is data scarcity: with only about twenty U.S.\ presidential
cycles (1948--2024), the shock-to-response mapping cannot be regressed directly
from history, so uncertainty must be modeled explicitly rather than assumed
away.

\section*{The approach and pipeline}
The system is a three-stage pipeline over three parallel demographic strata ---
race (5 blocs), religion (7), and gender (3), each independently covering the
whole electorate, deliberately without cross-tabulation.

\emph{Stage 1 --- language model.} A \texttt{Mistral-7B-v0.3} model, fine-tuned
with QLoRA (rank 16, $\alpha{=}32$, 4-bit NF4, learning rate $2\times10^{-4}$,
3 epochs), maps a shock description to per-bloc loyalty shifts across all 15
blocs. Output is emitted as one of nine ordinal ``delta bins''
(\texttt{strong\_neg} through \texttt{strong\_pos}) via constrained decoding
(the \texttt{outlines} library with a Pydantic schema), which guarantees every
token is schema-valid; the bins were originally decoded onto an ungrounded
$\pm0.15$ range and have since been rescaled to $\pm0.03$, anchored on the
largest single-shock per-bloc delta ever measured in the VOTER panel (see
``Magnitude calibration'' below). Held-out mean absolute error in delta units
was \textbf{0.0362} on the \emph{pre-corpus-rebalance Mistral v2} model
(111-example eval, under the OLD $\pm0.15$ decode scale)
\todo{re-run MAE on the current v3 model, against the 246-example eval set AND
the rescaled $\pm0.03$ decode table --- 0.0362 is stale on both the model
version and the scale, not comparable to a rescaled-scale figure}.

\emph{Stage 2 --- optimizer.} A disciplined quasi-convex program maximizes win
probability in Sharpe-ratio form,
$\max_w \Phi\!\big((\mu_{\text{eff}}(w)-V_{eq})/(\lambda_1\sqrt{w^\top\Sigma_\Delta
w})\big)$, where $\mu_{\text{eff}}$ is a $\lambda$-weighted blend of the three
strata and $V_{eq}$ is an electoral-college--adjusted win threshold
($0.5066$ Democrat, $0.4934$ Republican). The five race-bloc weights are the
decision variables, bounded to $[0.5\times,1.5\times]$ their baseline
population shares for demographic plausibility; the objective is provably
quasi-convex and solved as a second-order cone program. Minimum-variance
optimization is deliberately \emph{not} used --- it zeroes out small blocs, which
is mathematically valid but politically nonsensical. Because the gap between
blocs' baseline loyalties dominates any shock-scale perturbation to
$\mu_{\text{eff}}$, this solve is run once per party rather than once per
shock: the resulting \emph{equilibrium} coalition is empirically stable
against single-shock effects (``Headline contribution: coalition stability''
below), and a shock is evaluated against that fixed coalition rather than
being asked to move it.

\emph{Stage 3 --- Monte Carlo.} Small-sample uncertainty is propagated with a
Logistic-Normal isometric-log-ratio (ILR) simulation: coalition weights are
mapped to unconstrained space via a Helmert contrast, the race-level covariance
$\Sigma_\Delta$ is pushed through the ILR Jacobian, 10{,}000 draws are taken, and
each is mapped back through a softmax to a valid coalition. The result is a
win-probability point estimate with a 5th/95th-percentile confidence interval.
$\Sigma_\Delta$ is a race-only $5\times5$ matrix estimated by Ledoit--Wolf
shrinkage over \textbf{9 cycles (1992--2024, 8 first-differences)}; pre-1990
cycles are excluded because thin blocs (\texttt{asian}, \texttt{latino},
\texttt{other\_race}) had sparse or imputed sampling that more than doubled their
apparent variance.
% CONFLICT: an earlier decision described a 20-cycle (1948-2024) covariance;
% superseded 2026-06-29 by the 9-cycle 1990+ window used here.

\emph{Data.} Two distinct demographic inputs feed the model and must not be
conflated. Per-bloc \emph{loyalty} --- how Democratic each bloc votes --- is a
multi-source survey merge over NEP exit polls, ANES, GSS, and CES, each first
aggregated with its own sampling weights and then reconciled across sources by
inverse-standard-error weighting (Pew NPORS and the Democracy Fund VOTER panel are
wired into the loader but deferred), with a bloc-specific 1992 Perot two-party
correction. Electorate \emph{composition} --- how large each bloc is, used to form
the baseline coalition weights --- comes separately from NEP exit-poll shares. One
cycle is a known weak spot: ANES 2020's sampling-weight column is null, so it falls
back to equal-weighting (inflating the White Democratic share by $\sim$3.6\,pp) and
is down-weighted or excluded from the affected merges. Supervision signal also comes
from social media, tiered by how reliably a user's demographics can be inferred:
individual-bio posts feed both the mean and the covariance, while group- and
platform-level proxies feed the mean only, and pure language-fallback posts are
held out entirely --- a discipline that protects the optimizer's
mean/variance-independence assumption. Because historical shock--response labels
are scarce, the fine-tuning corpus is augmented with synthetic events generated
by DeepSeek and reviewed by Gemini over a balanced taxonomy of event type,
valence, electoral effect, and party. An earlier, label-skewed version (Mistral
v2 era) was rebalanced and the model retrained; the current merged corpus is
\textbf{1{,}228} records (\texttt{data/finetune/synthetic.jsonl}), split 982
train / 246 eval (\texttt{train.jsonl}, \texttt{eval.jsonl}), verified by line
count.

\section*{Headline contribution: neutrality validation}
The project's strongest result is a protocol for proving the system is
partisan-neutral --- and using it to find and fix a real bias. A
\textbf{zero-shock null test} feeds all-zero deltas through the full pipeline for
both parties; a neutral system must return symmetric baselines. It did not: the
Democratic baseline showed $\mu_{\text{eff}}=0.5575$ (margin $+0.0509$, win
probability $1.00$) against the Republican $0.5000$ (margin $+0.0066$, win
probability $0.26$). The cause was not the machine-learning model but an
\textbf{equal-within-stratum weighting placeholder} that over-weighted small,
strongly-Democratic religion and gender blocs by 7--33$\times$. Replacing it with
real electorate marginals cut the structural asymmetry from $+0.052$ to $+0.001$
--- a roughly \textbf{98\% reduction} --- with the tiny residual now slightly
favoring Republicans, exactly as the EC-adjusted threshold intends.

A \textbf{matched-pair symmetry test} (the same shock framed for each party) then
exposed a smaller residual behavioral lean of \textbf{$+0.028$} toward Democrats.
This was traced by elimination: baseline aggregation was fixed; the Monte Carlo
was switched to the real covariance; and the synthetic corpus was rebalanced to a
$-0.0007$ label asymmetry and the adapter retrained --- yet the lean persisted at
$+0.028$. The conclusion is that the residual comes from the base model's
pretrained political priors, which survive light LoRA fine-tuning and concentrate
in valence-interpretation events (e.g.\ a Republican gaffe read as more damaging
than its Democratic mirror). It is disclosed as a limitation rather than tuned
away, and a controlled comparison against a \texttt{Gemma-2-9b} model (identical
corpus and hyperparameters) is set up to test whether the lean is base-model
specific \todo{Gemma comparison run pending}.

\section*{Headline contribution: coalition stability}
A second empirical finding, stated positively: \textbf{under demographically
plausible coalition constraints, the optimal coalition is stable against
single-shock effects.} This is a substantive result about electoral
coalitions, not a limitation to disclose apologetically --- a campaign's
demographically bounded, risk-adjusted strategic emphasis does not rationally
re-litigate itself after every single news event. Bisection directly on the
optimizer found that coalition weights first move at a delta magnitude of
\textbf{$\sim$0.0633}, while the model's hard clip ceiling on any shock's
per-bloc delta is \textbf{0.0375} --- a full 1.7$\times$ below the movement
threshold, unreachable at any legal shock intensity. This was confirmed three
independent ways, none of which moved the equilibrium weights: adversarial
synthetic construction (every bloc pinned to its most extreme decode bin,
intensity multipliers swept to $10\times$); real model inference at the
API's maximum legal intensity ($3\times$); and both party framings (Democrat
and Republican baselines land at different corners of the
demographic-plausibility bound, but neither moves in response to a shock once
the party is fixed). Meanwhile $\mu_{\text{eff}}$ and win probability
\emph{do} respond to a shock --- under a tighter-margin baseline (Republican
framing), mean win-probability movement across tested shocks was $0.0503$,
with real, non-degenerate confidence intervals of roughly
$0.012$--$0.015$. The shock-conditional signal is real; it lives in
$\mu_{\text{eff}}$ and win probability, not in the coalition weights.
Widening the demographic-plausibility bounds until weights moved was
considered and rejected: that would mean choosing a bound because it makes
the output move, the same reasoning that originally produced the ungrounded
$\pm0.15$ delta-bin range below.

Read as a claim about American electoral coalitions rather than about this
model's arithmetic: \emph{coalition strategy is set by the structural,
slow-moving gap between how loyal different demographic blocs already are
to a party, not by the fast-moving news cycle of individual shocks.} Made
interpretable, $0.0633$ is $\sim$22$\times$ the panel's median measured
single-shock effect ($0.00287$), $\sim$14$\times$ the mean, and
$\sim$2.2$\times$ the single largest per-bloc shock effect ever recorded in
this project's ground truth ($0.02865$) --- moving the equilibrium coalition
would take an effect more than double anything actually observed, a
falsifiable comparison rather than a bare assertion. Only part of this claim
is unconditional, and the paper says so explicitly rather than letting the
stronger half borrow credibility from the weaker one: the loyalty gap
genuinely dwarfing any measured shock ($\sim$171$\times$ at the median,
$\sim$17$\times$ even at the largest shock on record) holds regardless of
the bounds; the exact $0.0633$ threshold does not, since it is a direct
function of the still-provisional $0.5$--$1.5\times$ bounds and would shift
if they did. The claim connects to, and partly tangles with, the literature
on partisan stability: it aligns with the social-identity account of
partisanship (Green, Palmquist \& Schickler 2002) and the Michigan-school
view of party identification as a stable filter for short-term information
(Campbell, Converse, Miller \& Stokes 1960), and with Key's (1955) account of
coalition change as concentrated in rare critical elections rather than
ordinary events. It sits in real tension with two findings: aggregate party
identification measured quarterly shows systematic short-term movement
(MacKuen, Erikson \& Stimson 1989), and individual events have been shown to
move vote share by amounts exceeding this project's own measured shocks ---
most strikingly Achen \& Bartels' (2016) finding that 1916 shark attacks cost
Woodrow Wilson up to ten points in coastal New Jersey, an effect larger than
anything in this paper's ground truth. A closer-magnitude example (Healy,
Malhotra \& Mo 2010, college-football wins moving vote share $\sim$1.6
points) was later shown likely not to replicate (Fowler \& Montagnes 2015).
The honest summary is that this project's finding matches the bulk of the
partisan-stability literature and the typical magnitude of the
retrospective-voting literature, while the shark-attack result shows ``no
single event moves a coalition this much'' is an empirical regularity in
this data, not a law of politics.

This is the same underlying fact as the magnitude-calibration result below,
seen from two angles: single-shock effects on bloc loyalty are small in
absolute terms (median $0.003$, largest ever measured $0.029$, once the
model's $3.44\times$ overstatement is corrected), and the optimal coalition
only responds to effects an order of magnitude larger than that. The
pre-rescale model got both halves of this story wrong in the same
direction --- overstating shock effects $\sim$4$\times$ pushed the optimizer
close enough to its bound-driven movement threshold to produce what looked
like genuine coalition responsiveness. Correcting the scale did not just fix
a calibration number; it resolved an apparent behavior that the corrected
numbers show was never really there.

\section*{Magnitude calibration}
The delta-bin decode range described above was, until recently, an
ungrounded $\pm0.15$ assumption, fixed before any ground-truth extraction
existed. Scored against the VOTER panel's magnitude-usable cells, the
model's predicted magnitudes were overstated by a mean cellwise
$|$predicted/measured$|$ ratio of \textbf{3.44$\times$}. The decode table was
rescaled to $\pm0.03$ (clip $\pm0.0375$), anchored on the largest
single-shock per-bloc delta ever measured in the panel (\textbf{0.0286}) ---
not a percentile, which would have made the largest real shocks in the
ground truth unrepresentable. Re-scored, the ratio falls to
\textbf{0.86$\times$}, close to unbiased. $\Sigma_\Delta$ was \emph{already}
estimated directly from real historical vote-share first-differences, in
real vote-share units, with no dependency on the decode table; the rescale
therefore did not introduce a new mismatch between the prediction scale and
the uncertainty model, it resolved a pre-existing one in which an overstated
prediction scale and an already-correctly-scaled uncertainty model had never
been checked against each other. Per-bloc \emph{direction} validation
remains a separate, harder problem: only 19 of 270 panel (shock, bloc) cells
pass the direction-trustworthiness filter (sign-stable, significant, not
suppressed, single-shock attributable), spanning 3 independent temporal
windows and 53\% concentrated in two blocs (\texttt{white},
\texttt{catholic}) --- a statistical-power failure, not a window-quality
one, and no ground-truth layer has any shock marked
\texttt{solely\_attributable} at all (survey layers measure windows spanning
multiple co-occurring events; the model predicts one event at a time). A
blinded expert-elicitation protocol is designed as the primary
direction-validation route (12 shocks, pre-registered Fleiss'-$\kappa$
agreement gate, magnitude buckets anchored to the rescaled $\pm0.03$ range)
but has not yet been run \todo{expert-elicitation results once the protocol
is executed}.

\section*{Honest limitations}
The three strata are combined additively without cross-tabulation, so the model
inherits the ecological fallacy (Robinson 1950); $\mu_{\text{eff}}$ and the
popular-vote threshold coincide only under that approximation. Only the five race
blocs are rebalanced, with religion and gender held fixed and treated as neutral
in the simulation, which understates uncertainty. The win threshold is a
model-internal effective-loyalty quantity, not a literal vote count. Synthetic
labels are LLM-generated, a documented deviation from the project's
reproducibility rule for this stage. And even after shrinkage the covariance
rests on only eight first-differences, so the thinnest blocs remain data-limited.
The residual $+0.028$ base-model lean is reported, not hidden.

\section*{Status}
The pipeline runs end to end and is deployed behind a streaming web interface.
The paper now anchors on two methodological contributions of equal weight: the
neutrality-validation finding, and the coalition-stability finding. Ground-truth
validation is split into two problems with different status: magnitude
calibration is solved and panel-anchored ($0.86\times$); per-bloc direction is
not, and depends on the expert-elicitation protocol, which is designed and
self-tested but not yet run. Remaining work is final-number consolidation
(fine-tuning MAE re-run against the rescaled decode table, the full symmetry
table, and the Gemma comparison), running the expert-elicitation protocol, IPF
raking to replace the fixed religion/gender weights, and a covariance
sensitivity analysis.
\todo{insert final results table once runs complete.}

\end{document}
